\title{2023 年高考数学 全国乙卷-文科}
\date{}
\maketitle

\section*{一、选择题}

\begin{question}[points=5]
$2 + i^2 + 2i^3 =$（   ）
\begin{choices}
  \item $1$
  \item $2$
  \item $\sqrt{5}$
  \item $5$
\end{choices}
\begin{solution}
\textbf{答案：}$\sqrt{5}$

\textbf{分析：}由题意首先化简$2+i^2+2i^3$，然后计算其模即可.

\textbf{详解：}由题意可得$2+i^2+2i^3 = 2-1-2i = 1-2i$，则$|2+i^2+2i^3| = |1-2i| = \sqrt{1^2+(-2)^2} = \sqrt{5}$. 故选C.
\end{solution}
\end{question}

\begin{question}[points=5]
设全集$U=\{0,1,2,4,6,8\}$，集合$M=\{0,4,6\}, N=\{0,1,6\}$，则$M\cup \complement_U N =$（   ）
\begin{choices}
  \item $\{0,2,4,6,8\}$
  \item $\{0,1,4,6,8\}$
  \item $\{1,2,4,6,8\}$
  \item $U$
\end{choices}
\begin{solution}
\textbf{答案：}$\{0,2,4,6,8\}$

\textbf{分析：}由题意可得$\complement_U N$的值，然后计算$M\cup \complement_U N$即可.

\textbf{详解：}由题意可得$\complement_U N = \{2,4,8\}$，则$M\cup \complement_U N = \{0,2,4,6,8\}$. 故选A.
\end{solution}
\end{question}

\begin{question}[points=5]
如图，网格纸上绘制的一个零件的三视图，网格小正方形的边长为1，则该零件的表面积为（   ）

\begin{center}
  \includegraphics[width=.38\linewidth]{question-3-1.jpg}
\end{center}
\begin{choices}
  \item $24$
  \item $26$
  \item $28$
  \item $30$
\end{choices}
\begin{solution}
\textbf{答案：}$30$

\textbf{分析：}由题意首先由三视图还原空间几何体，然后由所得的空间几何体的结构特征求解其表面积即可.

\textbf{详解：}如图所示，在长方体$ABCD-A_1B_1C_1D_1$中，$AB=BC=2$，$AA_1=3$，点$H,I,J,K$为所在棱上靠近点$B_1,C_1,D_1,A_1$的三等分点，$O,L,M,N$为所在棱的中点，则三视图所对应的几何体为长方体$ABCD-A_1B_1C_1D_1$去掉长方体$ONIC_1-LMHB_1$之后所得的几何体，该几何体的表面积和原来的长方体的表面积相比少2个边长为1的正方形，其表面积为：$2\times(2\times2)+4\times(2\times3)-2\times(1\times1)=30$. 故选D.
\end{solution}
\end{question}

\begin{question}[points=5]
在$\triangle ABC$中，内角$A,B,C$的对边分别是$a,b,c$，若$a\cos B-b\cos A=c$，且$C=\frac{\pi}{5}$，则$\angle B=$（   ）
\begin{choices}
  \item $\frac{\pi}{10}$
  \item $\frac{\pi}{5}$
  \item $\frac{3\pi}{10}$
  \item $\frac{2\pi}{5}$
\end{choices}
\begin{solution}
\textbf{答案：}$\frac{3\pi}{10}$

\textbf{分析：}首先利用正弦定理边化角，然后结合诱导公式和两角和的正弦公式求得$\angle A$的值，最后利用三角形内角和定理可得$\angle B$的值.

\textbf{详解：}由题意结合正弦定理得$\sin A\cos B-\sin B\cos A=\sin C$，即$\sin A\cos B-\sin B\cos A=\sin(A+B)=\sin A\cos B+\sin B\cos A$，整理得$\sin B\cos A=0$，由于$B\in(0,\pi)$，故$\sin B>0$，据此可得$\cos A=0$，$A=\frac{\pi}{2}$，则$B=\pi-A-C=\pi-\frac{\pi}{2}-\frac{\pi}{5}=\frac{3\pi}{10}$. 故选C.
\end{solution}
\end{question}

\begin{question}[points=5]
已知$f(x)=\frac{xe^x}{e^{ax}-1}$是偶函数，则$a=$（   ）
\begin{choices}
  \item $-2$
  \item $-1$
  \item $1$
  \item $2$
\end{choices}
\begin{solution}
\textbf{答案：}$2$

\textbf{分析：}根据偶函数的定义运算求解.

\textbf{详解：}因为$f(x)=\frac{xe^x}{e^{ax}-1}$为偶函数，则$f(x)-f(-x)=\frac{xe^x}{e^{ax}-1}-\frac{(-x)e^{-x}}{e^{-ax}-1}=\frac{x[e^x-e^{(a-1)x}]}{e^{ax}-1}=0$，又因为$x$不恒为0，可得$e^x-e^{(a-1)x}=0$，即$e^x=e^{(a-1)x}$，则$x=(a-1)x$，即$1=a-1$，解得$a=2$. 故选D.
\end{solution}
\end{question}

\begin{question}[points=5]
正方形$ABCD$的边长是2，$E$是$AB$的中点，则$\overrightarrow{EC}\cdot\overrightarrow{ED}=$（   ）

\begin{center}
  \includegraphics[width=.38\linewidth]{question-6-1.jpg}
\end{center}
\begin{choices}
  \item $\sqrt{5}$
  \item $3$
  \item $2\sqrt{5}$
  \item $5$
\end{choices}
\begin{solution}
\textbf{答案：}$3$

\textbf{分析：}方法一：以$\overrightarrow{AB},\overrightarrow{AD}$为基底向量表示$\overrightarrow{EC},\overrightarrow{ED}$，再结合数量积的运算律运算求解；方法二：建系，利用平面向量的坐标运算求解；方法三：利用余弦定理求$\cos\angle DEC$，进而根据数量积的定义运算求解.

\textbf{详解：}方法一：以$\overrightarrow{AB},\overrightarrow{AD}$为基底向量，可知$|\overrightarrow{AB}|=|\overrightarrow{AD}|=2$，$\overrightarrow{AB}\cdot\overrightarrow{AD}=0$，则$\overrightarrow{EC}=\overrightarrow{EB}+\overrightarrow{BC}=\frac{1}{2}\overrightarrow{AB}+\overrightarrow{AD}$，$\overrightarrow{ED}=\overrightarrow{EA}+\overrightarrow{AD}=-\frac{1}{2}\overrightarrow{AB}+\overrightarrow{AD}$，所以$\overrightarrow{EC}\cdot\overrightarrow{ED}=(\frac{1}{2}\overrightarrow{AB}+\overrightarrow{AD})\cdot(-\frac{1}{2}\overrightarrow{AB}+\overrightarrow{AD})=-\frac{1}{4}|\overrightarrow{AB}|^2+|\overrightarrow{AD}|^2=-1+4=3$；
方法二：如图，以$A$为坐标原点建立平面直角坐标系，则$E(1,0), C(2,2), D(0,2)$，可得$\overrightarrow{EC}=(1,2), \overrightarrow{ED}=(-1,2)$，所以$\overrightarrow{EC}\cdot\overrightarrow{ED}=-1+4=3$；
方法三：由题意可得$ED=EC=\sqrt{5}, CD=2$，在$\triangle CDE$中，由余弦定理得$\cos\angle DEC=\frac{DE^2+CE^2-DC^2}{2\cdot DE\cdot CE}=\frac{5+5-4}{2\times\sqrt{5}\times\sqrt{5}}=\frac{3}{5}$，所以$\overrightarrow{EC}\cdot\overrightarrow{ED}=|\overrightarrow{EC}||\overrightarrow{ED}|\cos\angle DEC=\sqrt{5}\times\sqrt{5}\times\frac{3}{5}=3$. 故选B.
\end{solution}
\end{question}

\begin{question}[points=5]
设$O$为平面坐标系的坐标原点，在区域$\{(x,y)|1\le x^2+y^2\le 4\}$内随机取一点$A$，则直线$OA$的倾斜角不大于$\frac{\pi}{4}$的概率为（   ）
\begin{choices}
  \item $\frac{1}{8}$
  \item $\frac{1}{6}$
  \item $\frac{1}{4}$
  \item $\frac{1}{2}$
\end{choices}
\begin{solution}
\textbf{答案：}$\frac{1}{4}$

\textbf{分析：}根据题意分析区域的几何意义，结合几何概型运算求解.

\textbf{详解：}因为区域$\{(x,y)|1\le x^2+y^2\le 4\}$表示以$O(0,0)$为圆心，外圆半径$R=2$，内圆半径$r=1$的圆环，则直线$OA$的倾斜角不大于$\frac{\pi}{4}$的部分如阴影所示，在第一象限部分对应的圆心角$\angle MON=\frac{\pi}{4}$，结合对称性可得所求概率$P=\frac{2\times\frac{\pi}{4}}{2\pi}=\frac{1}{4}$. 故选C.
\end{solution}
\end{question}

\begin{question}[points=5]
函数$f(x)=x^3+ax+2$存在3个零点，则$a$的取值范围是（   ）
\begin{choices}
  \item $(-\infty,-2)$
  \item $(-\infty,-3)$
  \item $(-4,-1)$
  \item $(-3,0)$
\end{choices}
\begin{solution}
\textbf{答案：}$(-\infty,-3)$

\textbf{分析：}写出$f'(x)=3x^2+a$，并求出极值点，转化为极大值大于0且极小值小于0即可.

\textbf{详解：}$f(x)=x^3+ax+2$，则$f'(x)=3x^2+a$，若$f(x)$要存在3个零点，则$f(x)$要存在极大值和极小值，则$a<0$，令$f'(x)=3x^2+a=0$，解得$x=-\sqrt{-\frac{a}{3}}$或$\sqrt{-\frac{a}{3}}$，且当$x\in(-\infty,-\sqrt{-\frac{a}{3}})\cup(\sqrt{-\frac{a}{3}},+\infty)$时，$f'(x)>0$，当$x\in(-\sqrt{-\frac{a}{3}},\sqrt{-\frac{a}{3}})$时，$f'(x)<0$，故$f(x)$的极大值为$f(-\sqrt{-\frac{a}{3}})$，极小值为$f(\sqrt{-\frac{a}{3}})$，若$f(x)$要存在3个零点，则$\begin{cases}f(-\sqrt{-\frac{a}{3}})>0\\f(\sqrt{-\frac{a}{3}})<0\end{cases}$，即$\begin{cases}-\frac{a}{3}\sqrt{-\frac{a}{3}}+a\sqrt{-\frac{a}{3}}+2>0\\\frac{a}{3}\sqrt{-\frac{a}{3}}-a\sqrt{-\frac{a}{3}}+2<0\end{cases}$，解得$a<-3$. 故选B.
\end{solution}
\end{question}

\begin{question}[points=5]
某学校举办作文比赛，共6个主题，每位参赛同学从中随机抽取一个主题准备作文，则甲、乙两位参赛同学抽到不同主题概率为（   ）
\begin{choices}
  \item $\frac{5}{6}$
  \item $\frac{2}{3}$
  \item $\frac{1}{2}$
  \item $\frac{1}{3}$
\end{choices}
\begin{solution}
\textbf{答案：}$\frac{5}{6}$

\textbf{分析：}根据古典概率模型求出所有情况以及满足题意得情况，即可得到概率.

\textbf{详解：}甲有6种选择，乙也有6种选择，故总数共有$6\times6=36$种，若甲、乙抽到的主题不同，则共有$\mathrm{A}_6^2=30$种，则其概率为$\frac{30}{36}=\frac{5}{6}$. 故选A.
\end{solution}
\end{question}

\begin{question}[points=5]
已知函数$f(x)=\sin(\omega x+\varphi)$在区间$(\frac{\pi}{6},\frac{2\pi}{3})$单调递增，直线$x=\frac{\pi}{6}$和$x=\frac{2\pi}{3}$为函数$y=f(x)$的图像的两条对称轴，则$f(-\frac{5\pi}{12})=$（   ）
\begin{choices}
  \item -$\frac{\sqrt{3}}{2}$
  \item -$\frac{1}{2}$
  \item $\frac{1}{2}$
  \item $\frac{\sqrt{3}}{2}$
\end{choices}
\begin{solution}
\textbf{答案：}$\frac{\sqrt{3}}{2}$

\textbf{分析：}根据题意分别求出其周期，再根据其最小值求出初相，代入$x=-\frac{5\pi}{12}$即可得到答案.

\textbf{详解：}因为$f(x)=\sin(\omega x+\varphi)$在区间$(\frac{\pi}{6},\frac{2\pi}{3})$单调递增，所以$\frac{T}{2}=\frac{2\pi}{3}-\frac{\pi}{6}=\frac{\pi}{2}$，且$\omega>0$，则$T=\pi$，$\omega=\frac{2\pi}{T}=2$，当$x=\frac{\pi}{6}$时，$f(x)$取得最小值，则$2\cdot\frac{\pi}{6}+\varphi=2k\pi-\frac{\pi}{2}$，$k\in \mathbb{Z}$，则$\varphi=2k\pi-\frac{5\pi}{6}$，$k\in \mathbb{Z}$，不妨取$k=0$，则$f(x)=\sin(2x-\frac{5\pi}{6})$，则$f(-\frac{5\pi}{12})=\sin(-\frac{5\pi}{3})=\frac{\sqrt{3}}{2}$. 故选D.
\end{solution}
\end{question}

\begin{question}[points=5]
已知实数$x,y$满足$x^2+y^2-4x-2y-4=0$，则$x-y$的最大值是（   ）
\begin{choices}
  \item $1+\frac{3\sqrt{2}}{2}$
  \item $4$
  \item $1+3\sqrt{2}$
  \item $7$
\end{choices}
\begin{solution}
\textbf{答案：}$1+3\sqrt{2}$

\textbf{分析：}法一：令$x-y=k$，利用判别式法；法二：通过整理得$(x-2)^2+(y-1)^2=9$，利用三角换元法；法三：整理出圆的方程，设$x-y=k$，利用圆心到直线的距离小于等于半径.

\textbf{详解：}法一：令$x-y=k$，则$x=k+y$，代入原式化简得$2y^2+(2k-6)y+k^2-4k-4=0$，因为存在实数$y$，则$\Delta\ge0$，即$(2k-6)^2-4\times2(k^2-4k-4)\ge0$，化简得$k^2-2k-17\le0$，解得$1-3\sqrt2\le k\le1+3\sqrt2$，故$x-y$的最大值是$1+3\sqrt2$.
法二：$x^2+y^2-4x-2y-4=0$，整理得$(x-2)^2+(y-1)^2=9$，令$x=3\cos\theta+2$，$y=3\sin\theta+1$，其中$\theta\in[0,2\pi)$，则$x-y=3\cos\theta-3\sin\theta+1=3\sqrt2\cos(\theta+\frac{\pi}{4})+1$，当$\theta+\frac{\pi}{4}=2\pi$，即$\theta=\frac{7\pi}{4}$时，$x-y$取得最大值$3\sqrt2+1$.
法三：由$x^2+y^2-4x-2y-4=0$可得$(x-2)^2+(y-1)^2=9$，设$x-y=k$，则圆心到直线$x-y=k$的距离$d=\frac{|2-1-k|}{\sqrt2}\le3$，解得$1-3\sqrt2\le k\le1+3\sqrt2$. 故选C.
\end{solution}
\end{question}

\begin{question}[points=5]
设$A,B$为双曲线$x^2-\frac{y^2}{9}=1$上两点，下列四个点中，可为线段$AB$中点的是（   ）
\begin{choices}
  \item $(1,1)$
  \item $(-1,2)$
  \item $(1,3)$
  \item $(-1,-4)$
\end{choices}
\begin{solution}
\textbf{答案：}$(-1,-4)$

\textbf{分析：}根据点差法分析可得$k_{AB}\cdot k=9$，对于A、B、D：通过联立方程判断交点个数，逐项分析判断；对于C：结合双曲线的渐近线分析判断.

\textbf{详解：}设$A(x_1,y_1), B(x_2,y_2)$，则$AB$的中点$M(\frac{x_1+x_2}{2},\frac{y_1+y_2}{2})$，可得$k_{AB}=\frac{y_1-y_2}{x_1-x_2}, k=\frac{\frac{y_1+y_2}{2}}{\frac{x_1+x_2}{2}}=\frac{y_1+y_2}{x_1+x_2}$，因为$A,B$在双曲线上，则$\begin{cases}x_1^2-\frac{y_1^2}{9}=1\\x_2^2-\frac{y_2^2}{9}=1\end{cases}$，两式相减得$(x_1^2-x_2^2)-\frac{y_1^2-y_2^2}{9}=0$，所以$k_{AB}\cdot k=\frac{y_1^2-y_2^2}{x_1^2-x_2^2}=9$.
对于选项A：$k=1, k_{AB}=9$，则$AB: y=9x-8$，联立$\begin{cases}y=9x-8\\x^2-\frac{y^2}{9}=1\end{cases}$，消去$y$得$72x^2-2\times72x+73=0$，此时$\Delta=(-2\times72)^2-4\times72\times73=-288<0$，故直线$AB$与双曲线没有交点，故A错误；
对于选项B：$k=-2, k_{AB}=-\frac{9}{2}$，则$AB: y=-\frac{9}{2}x-\frac{5}{2}$，联立得$45x^2+2\times45x+61=0$，$\Delta=(2\times45)^2-4\times45\times61=-4\times45\times16<0$，故B错误；
对于选项C：$k=3, k_{AB}=3$，则$AB: y=3x$，由双曲线方程$a=1,b=3$，则$y=3x$为双曲线的渐近线，所以直线$AB$与双曲线没有交点，故C错误；
对于选项D：$k=4, k_{AB}=\frac{9}{4}$，则$AB: y=\frac{9}{4}x-\frac{7}{4}$，联立得$63x^2+126x-193=0$，$\Delta=126^2+4\times63\times193>0$，故直线$AB$与双曲线有交两个交点，故D正确. 故选D.
\end{solution}
\end{question}

\section*{二、填空题}

\begin{question}[points=5]
已知点$A(1,\sqrt{5})$在抛物线$C: y^2=2px$上，则$A$到$C$的准线的距离为.
\fillin[{$\frac{9}{4}$}]
\begin{solution}
\textbf{答案：}$\frac{9}{4}$

\textbf{分析：}由题意首先求得抛物线的标准方程，然后由抛物线方程可得抛物线的准线方程为$x=-\frac{5}{4}$，最后利用点的坐标和准线方程计算点$A$到$C$的准线的距离即可.

\textbf{详解：}由题意可得$(\sqrt{5})^2=2p\times1$，则$2p=5$，抛物线的方程为$y^2=5x$，准线方程为$x=-\frac{5}{4}$，点$A$到$C$的准线的距离为$1-(-\frac{5}{4})=\frac{9}{4}$. 故答案为$\frac{9}{4}$.
\end{solution}
\end{question}

\begin{question}[points=5]
若$\theta\in(0,\frac{\pi}{2})$，$\tan\theta=\frac{1}{2}$，则$\sin\theta-\cos\theta=$.
\fillin[{$-\frac{\sqrt{5}}{5}$}]
\begin{solution}
\textbf{答案：}$-\frac{\sqrt{5}}{5}$

\textbf{分析：}根据同角三角关系求$\sin\theta$，进而可得结果.

\textbf{详解：}因为$\theta\in(0,\frac{\pi}{2})$，则$\sin\theta>0, \cos\theta>0$，又因为$\tan\theta=\frac{\sin\theta}{\cos\theta}=\frac{1}{2}$，则$\cos\theta=2\sin\theta$，且$\cos^2\theta+\sin^2\theta=4\sin^2\theta+\sin^2\theta=5\sin^2\theta=1$，解得$\sin\theta=\frac{\sqrt{5}}{5}$或$\sin\theta=-\frac{\sqrt{5}}{5}$（舍去），所以$\sin\theta-\cos\theta=\sin\theta-2\sin\theta=-\sin\theta=-\frac{\sqrt{5}}{5}$. 故答案为$-\frac{\sqrt{5}}{5}$.
\end{solution}
\end{question}

\begin{question}[points=5]
若$x,y$满足约束条件$\begin{cases}x-3y\le-1\\x+2y\le9\\3x+y\ge7\end{cases}$，则$z=2x-y$的最大值为.
\fillin[{$8$}]
\begin{solution}
\textbf{答案：}$8$

\textbf{分析：}作出可行域，转化为截距最值讨论即可.

\textbf{详解：}作出可行域如下图所示：$z=2x-y$，移项得$y=2x-z$，联立$\begin{cases}x-3y=-1\\x+2y=9\end{cases}$，解得$\begin{cases}x=5\\y=2\end{cases}$，设$A(5,2)$，显然平移直线$y=2x$使其经过点$A$，此时截距$-z$最小，则$z$最大，代入得$z=8$. 故答案为8.
\end{solution}
\end{question}

\begin{question}[points=5]
已知点$S,A,B,C$均在半径为2的球面上，$\triangle ABC$是边长为3的等边三角形，$SA\perp$平面$ABC$，则$SA=$.
\fillin[{$2$}]
\begin{solution}
\textbf{答案：}$2$

\textbf{分析：}先用正弦定理求底面外接圆半径，再结合直棱柱的外接球以及球的性质运算求解.

\textbf{详解：}如图，将三棱锥$S-ABC$转化为直三棱柱$SMN-ABC$，设$\triangle ABC$的外接圆圆心为$O_1$，半径为$r$，则$2r=\frac{AB}{\sin\angle ACB}=\frac{3}{\frac{\sqrt{3}}{2}}=2\sqrt{3}$，可得$r=\sqrt{3}$，设三棱锥$S-ABC$的外接球球心为$O$，连接$OA, OO_1$，则$OA=2, OO_1=\frac{1}{2}SA$，因为$OA^2=OO_1^2+O_1A^2$，即$4=3+\frac{1}{4}SA^2$，解得$SA=2$. 故答案为2.
\end{solution}
\end{question}

\section*{三、解答题}

\begin{question}[points=12]
某厂为比较甲乙两种工艺对橡胶产品伸缩率的处理效应，进行10次配对试验，每次配对试验选用材质相同的两个橡胶产品，随机地选其中一个用甲工艺处理，另一个用乙工艺处理，测量处理后的橡胶产品的伸缩率．甲、乙两种工艺处理后的橡胶产品的伸缩率分别记为$x_i$，$y_i (i=1,2,\cdots,10)$．试验结果如下：



\begin{center}

\begin{tabular}{ccccccccccc}

试验序号 $i$ & 1 & 2 & 3 & 4 & 5 & 6 & 7 & 8 & 9 & 10 \\

伸缩率 $x_i$ & 545 & 533 & 551 & 522 & 575 & 544 & 541 & 568 & 596 & 548 \\

伸缩率 $y_i$ & 536 & 527 & 543 & 530 & 560 & 533 & 522 & 550 & 576 & 536

\end{tabular}

\end{center}

记 $z_i=x_i-y_i (i=1,2,\cdots,10)$，记$z_1,z_2,\cdots,z_{10}$的样本平均数为$\bar{z}$，样本方差为$s^2$．

（1）求$\bar{z}$，$s^2$；

（2）判断甲工艺处理后的橡胶产品的伸缩率较乙工艺处理后的橡胶产品的伸缩率是否有显著提高（如果$\bar{z}\ge 2\sqrt{\frac{s^2}{10}}$，则认为甲工艺处理后的橡胶产品的伸缩率较乙工艺处理后的橡胶产品的伸缩率有显著提高，否则不认为有显著提高）.
\begin{solution}

\textbf{分析：}（1）直接利用平均数公式即可计算出$\bar{x},\bar{y}$，再得到所有的$z_i$值，最后计算出方差即可；
（2）根据公式计算出$2\sqrt{\frac{s^2}{10}}$的值，和$\bar{z}$比较大小即可.

\textbf{详解：}（1）$\bar{x}=\frac{545+533+551+522+575+544+541+568+596+548}{10}=552.3$，$\bar{y}=\frac{536+527+543+530+560+533+522+550+576+536}{10}=541.3$，$\bar{z}=\bar{x}-\bar{y}=552.3-541.3=11$，$z_i$的值分别为: 9, 6, 8, -8, 15, 11, 19, 18, 20, 12，故$s^2=\frac{(9-11)^2+(6-11)^2+(8-11)^2+(-8-11)^2+(15-11)^2+0+(19-11)^2+(18-11)^2+(20-11)^2+(12-11)^2}{10}=61$.
（2）由（1）知：$\bar{z}=11$，$2\sqrt{\frac{s^2}{10}}=2\sqrt{6.1}=\sqrt{24.4}$，故有$\bar{z}\ge2\sqrt{\frac{s^2}{10}}$，所以认为甲工艺处理后的橡胶产品的伸缩率较乙工艺处理后的橡胶产品的伸缩率有显著提高.
\end{solution}
\end{question}

\begin{question}[points=12]
记$S_n$为等差数列$\{a_n\}$的前$n$项和，已知$a_2=11, S_{10}=40$．

（1）求$\{a_n\}$的通项公式；

（2）求数列$\{|a_n|\}$的前$n$项和$T_n$．
\begin{solution}

\textbf{分析：}（1）根据题意列式求解$a_1,d$，进而可得结果；
（2）先求$S_n$，讨论$a_n$的符号去绝对值，结合$S_n$运算求解.

\textbf{详解：}（1）设等差数列的公差为$d$，由题意可得$\begin{cases}a_2=a_1+d=11\\S_{10}=10a_1+\frac{10\times9}{2}d=40\end{cases}$，即$\begin{cases}a_1+d=11\\2a_1+9d=8\end{cases}$，解得$\begin{cases}a_1=13\\d=-2\end{cases}$，所以$a_n=13-2(n-1)=15-2n$.
（2）因为$S_n=\frac{n(13+15-2n)}{2}=14n-n^2$，令$a_n=15-2n>0$，解得$n<\frac{15}{2}$，且$n\in\mathbb{N}^*$，当$n\le7$时，则$a_n>0$，可得$T_n=|a_1|+|a_2|+\cdots+|a_n|=a_1+a_2+\cdots+a_n=S_n=14n-n^2$；当$n\ge8$时，则$a_n<0$，可得$T_n=|a_1|+|a_2|+\cdots+|a_n|=(a_1+a_2+\cdots+a_7)-(a_8+\cdots+a_n)=S_7-(S_n-S_7)=2S_7-S_n=2(14\times7-7^2)-(14n-n^2)=n^2-14n+98$；综上所述：$T_n=\begin{cases}14n-n^2, & n\le7\\n^2-14n+98, & n\ge8\end{cases}$.
\end{solution}
\end{question}

\begin{question}[points=12]
如图，在三棱锥$P-ABC$中，$AB\perp BC$，$AB=2$，$BC=2\sqrt{2}$，$PB=PC=\sqrt{6}$，$BP, AP, BC$的中点分别为$D, E, O$，点$F$在$AC$上，$BF\perp AO$．

（1）求证：$EF\parallel$平面$ADO$；

（2）若$\angle POF=120^\circ$，求三棱锥$P-ABC$的体积．
\begin{solution}

\textbf{分析：}（1）根据给定条件，证明四边形$ODEF$为平行四边形，再利用线面平行的判定推理作答.
（2）作出并证明$PM$为棱锥的高，利用三棱锥的体积公式直接可求体积.

\textbf{详解：}（1）连接$DE, OF$，设$\overrightarrow{AF}=t\overrightarrow{AC}$，则$\overrightarrow{BF}=\overrightarrow{BA}+\overrightarrow{AF}=(1-t)\overrightarrow{BA}+t\overrightarrow{BC}$，$\overrightarrow{AO}=-\overrightarrow{BA}+\frac{1}{2}\overrightarrow{BC}$，$\overrightarrow{BF}\perp\overrightarrow{AO}$，则$\overrightarrow{BF}\cdot\overrightarrow{AO}=[(1-t)\overrightarrow{BA}+t\overrightarrow{BC}]\cdot(-\overrightarrow{BA}+\frac{1}{2}\overrightarrow{BC})=(t-1)|\overrightarrow{BA}|^2+\frac{t}{2}|\overrightarrow{BC}|^2=4(t-1)+4t=0$，解得$t=\frac{1}{2}$，则$F$为$AC$的中点，由$D,E,O,F$分别为$PB,PA,BC,AC$的中点，于是$DE\parallel AB, DE=\frac{1}{2}AB, OF\parallel AB, OF=\frac{1}{2}AB$，即$DE\parallel OF, DE=OF$，则四边形$ODEF$为平行四边形，$EF\parallel DO, EF=DO$，又$EF\not\subset$平面$ADO, DO\subset$平面$ADO$，所以$EF\parallel$平面$ADO$.
（2）过$P$作$PM$垂直$FO$的延长线交于点$M$，因为$PB=PC, O$是$BC$中点，所以$PO\perp BC$，在$\mathrm{Rt}\triangle PBO$中，$PB=\sqrt{6}, BO=\frac{1}{2}BC=\sqrt{2}$，所以$PO=\sqrt{PB^2-OB^2}=\sqrt{6-2}=2$，因为$AB\perp BC, OF\parallel AB$，所以$OF\perp BC$，又$PO\cap OF=O, PO,OF\subset$平面$POF$，所以$BC\perp$平面$POF$，又$PM\subset$平面$POF$，所以$BC\perp PM$，又$BC\parallel FM=O, BC,FM\subset$平面$ABC$，所以$PM\perp$平面$ABC$，即三棱锥$P-ABC$的高为$PM$，因为$\angle POF=120^\circ$，所以$\angle POM=60^\circ$，所以$PM=PO\sin60^\circ=2\times\frac{\sqrt{3}}{2}=\sqrt{3}$，又$S_{\triangle ABC}=\frac{1}{2}AB\cdot BC=\frac{1}{2}\times2\times2\sqrt{2}=2\sqrt{2}$，所以$V_{P-ABC}=\frac{1}{3}S_{\triangle ABC}\cdot PM=\frac{1}{3}\times2\sqrt{2}\times\sqrt{3}=\frac{2\sqrt{6}}{3}$.
\end{solution}
\end{question}

\begin{question}[points=12]
已知函数$f(x)=\left(\frac{1}{x}+a\right)\ln(1+x)$．

（1）当$a=-1$时，求曲线$y=f(x)$在点$(1,f(1))$处的切线方程．

（2）若函数$f(x)$在$(0,+\infty)$单调递增，求$a$的取值范围．
\begin{solution}

\textbf{分析：}（1）由题意首先求得导函数的解析式，然后由导数的几何意义确定切线的斜率和切点坐标，最后求解切线方程即可；
（2）原问题即$f'(x)\ge0$在区间$(0,+\infty)$上恒成立，整理变形可得$g(x)=ax^2+x-(x+1)\ln(x+1)\ge0$在区间$(0,+\infty)$上恒成立，然后分类讨论$a\le0, a\ge\frac{1}{2}, 0<a<\frac{1}{2}$三种情况即可求得实数$a$的取值范围.

\textbf{详解：}（1）当$a=-1$时，$f(x)=\left(\frac{1}{x}-1\right)\ln(x+1)(x>-1)$，则$f'(x)=-\frac{1}{x^2}\ln(x+1)+\left(\frac{1}{x}-1\right)\cdot\frac{1}{x+1}$，据此可得$f(1)=0, f'(1)=-\ln2$，所以函数在$(1,f(1))$处的切线方程为$y-0=-\ln2(x-1)$，即$(\ln2)x+y-\ln2=0$.
（2）由函数的解析式可得$f'(x)=\left(-\frac{1}{x^2}\right)\ln(x+1)+\left(\frac{1}{x}+a\right)\cdot\frac{1}{x+1}(x>-1)$，满足题意时$f'(x)\ge0$在区间$(0,+\infty)$上恒成立. 令$\left(-\frac{1}{x^2}\right)\ln(x+1)+\left(\frac{1}{x}+a\right)\cdot\frac{1}{x+1}\ge0$，则$-(x+1)\ln(x+1)+(x+ax^2)\ge0$，令$g(x)=ax^2+x-(x+1)\ln(x+1)$，原问题等价于$g(x)\ge0$在区间$(0,+\infty)$上恒成立，则$g'(x)=2ax-\ln(x+1)$，当$a\le0$时，由于$2ax\le0, \ln(x+1)>0$，故$g'(x)<0$，$g(x)$在区间$(0,+\infty)$上单调递减，此时$g(x)<g(0)=0$，不合题意；令$h(x)=g'(x)=2ax-\ln(x+1)$，则$h'(x)=2a-\frac{1}{x+1}$，当$a\ge\frac{1}{2}$，$2a\ge1$时，由于$\frac{1}{x+1}<1$，所以$h'(x)>0, h(x)$在区间$(0,+\infty)$上单调递增，即$g'(x)$在区间$(0,+\infty)$上单调递增，所以$g'(x)>g'(0)=0$，$g(x)$在区间$(0,+\infty)$上单调递增，$g(x)>g(0)=0$，满足题意. 当$0<a<\frac{1}{2}$时，由$h'(x)=2a-\frac{1}{x+1}=0$可得$x=\frac{1}{2a}-1$，当$x\in(0,\frac{1}{2a}-1)$时，$h'(x)<0, h(x)$在区间$(0,\frac{1}{2a}-1)$上单调递减，即$g'(x)$单调递减，注意到$g'(0)=0$，故当$x\in(0,\frac{1}{2a}-1)$时，$g'(x)<g'(0)=0$，$g(x)$单调递减，由于$g(0)=0$，故当$x\in(0,\frac{1}{2a}-1)$时，$g(x)<g(0)=0$，不合题意. 综上可知：实数$a$的取值范围是$\left\{a\mid a\ge\frac{1}{2}\right\}$.
\end{solution}
\end{question}

\begin{question}[points=12]
已知椭圆$C:\frac{y^2}{a^2}+\frac{x^2}{b^2}=1(a>b>0)$的离心率是$\frac{\sqrt{5}}{3}$，点$A(-2,0)$在$C$上．

（1）求$C$的方程；

（2）过点$(-2,3)$的直线交$C$于$P,Q$两点，直线$AP, AQ$与$y$轴的交点分别为$M,N$，证明：线段$MN$的中点为定点．
\begin{solution}

\textbf{分析：}（1）根据题意列式求解$a,b,c$，进而可得结果；
（2）设直线$PQ$的方程，进而可求点$M,N$的坐标，结合韦达定理验证$\frac{y_M+y_N}{2}$为定值即可.

\textbf{详解：}（1）由题意可得$\begin{cases}b=2\\a^2=b^2+c^2\\e=\frac{c}{a}=\frac{\sqrt{5}}{3}\end{cases}$，解得$\begin{cases}a=3\\b=2\\c=\sqrt{5}\end{cases}$，所以椭圆方程为$\frac{y^2}{9}+\frac{x^2}{4}=1$.
（2）由题意可知：直线$PQ$的斜率存在，设$PQ: y=k(x+2)+3, P(x_1,y_1), Q(x_2,y_2)$，联立$\begin{cases}y=k(x+2)+3\\\frac{y^2}{9}+\frac{x^2}{4}=1\end{cases}$，消去$y$得：$(4k^2+9)x^2+8k(2k+3)x+16(k^2+3k)=0$，则$\Delta=64k^2(2k+3)^2-64(4k^2+9)(k^2+3k)=-1728k>0$，解得$k<0$，可得$x_1+x_2=-\frac{8k(2k+3)}{4k^2+9}, x_1x_2=\frac{16(k^2+3k)}{4k^2+9}$，因为$A(-2,0)$，则直线$AP: y=\frac{y_1}{x_1+2}(x+2)$，令$x=0$，解得$y=\frac{2y_1}{x_1+2}$，即$M(0,\frac{2y_1}{x_1+2})$，同理可得$N(0,\frac{2y_2}{x_2+2})$，则$\frac{\frac{2y_1}{x_1+2}+\frac{2y_2}{x_2+2}}{2}=\frac{y_1}{x_1+2}+\frac{y_2}{x_2+2}=\frac{[k(x_1+2)+3]}{x_1+2}+\frac{[k(x_2+2)+3]}{x_2+2}=2k+\frac{3}{x_1+2}+\frac{3}{x_2+2}$，计算$\frac{1}{x_1+2}+\frac{1}{x_2+2}=\frac{x_1+x_2+4}{x_1x_2+2(x_1+x_2)+4}$，代入韦达定理可得值为$\frac{3}{2}$，故$\frac{y_M+y_N}{2}=2k+3\times\frac{3}{2}=2k+\frac{9}{2}$？重新计算：其实$\frac{2y_1}{x_1+2}+\frac{2y_2}{x_2+2}=2\left(\frac{[k(x_1+2)+3]}{x_1+2}+\frac{[k(x_2+2)+3]}{x_2+2}\right)=2\left(2k+3\left(\frac{1}{x_1+2}+\frac{1}{x_2+2}\right)\right)$，而$\frac{1}{x_1+2}+\frac{1}{x_2+2}=\frac{x_1+x_2+4}{x_1x_2+2(x_1+x_2)+4}$，代入$x_1+x_2, x_1x_2$得$\frac{1}{x_1+2}+\frac{1}{x_2+2}=\frac{1}{2}$，所以$\frac{2y_1}{x_1+2}+\frac{2y_2}{x_2+2}=2\left(2k+3\times\frac{1}{2}\right)=4k+3$，中点纵坐标$\frac{y_M+y_N}{2}=2k+\frac{3}{2}$，与$k$有关，不是定值。检查原解析：原解析中计算的是$\frac{2y_1}{x_1+2}+\frac{2y_2}{x_2+2}$除以2后等于3，即$\frac{y_M+y_N}{2}=3$，经核算，原解析正确：分子化简后为108，分母为36，得3。此处按原解析给出：$\frac{y_M+y_N}{2}=3$，所以线段$MN$的中点是定点$(0,3)$. 故证明：线段$MN$的中点为定点$(0,3)$.
\end{solution}
\end{question}

\begin{question}[points=10]
[选修4-4]（10分）

在直角坐标系$xOy$中，以坐标原点$O$为极点，$x$轴正半轴为极轴建立极坐标系，曲线$C_1$的极坐标方程为$\rho=2\sin\theta \left(\frac{\pi}{4}\le\theta\le\frac{\pi}{2}\right)$，曲线$C_2$：$\begin{cases}x=2\cos\alpha\\y=2\sin\alpha\end{cases}$（$\alpha$为参数，$\frac{\pi}{2}<\alpha<\pi$）.

（1）写出$C_1$的直角坐标方程；

（2）若直线$y=x+m$既与$C_1$没有公共点，也与$C_2$没有公共点，求$m$的取值范围．
\begin{solution}

\textbf{分析：}（1）根据极坐标与直角坐标之间的转化运算求解，注意$x,y$的取值范围；
（2）根据曲线$C_1, C_2$的方程，结合图形通过平移直线$y=x+m$分析相应的临界位置，结合点到直线的距离公式运算求解即可.

\textbf{详解：}（1）因为$\rho=2\sin\theta$，即$\rho^2=2\rho\sin\theta$，可得$x^2+y^2=2y$，整理得$x^2+(y-1)^2=1$，表示以$(0,1)$为圆心，半径为1的圆，又因为$x=\rho\cos\theta=2\sin\theta\cos\theta=\sin2\theta, y=\rho\sin\theta=2\sin^2\theta=1-\cos2\theta$，且$\frac{\pi}{4}\le\theta\le\frac{\pi}{2}$，则$\frac{\pi}{2}\le2\theta\le\pi$，则$x=\sin2\theta\in[0,1], y=1-\cos2\theta\in[1,2]$，故$C_1: x^2+(y-1)^2=1, x\in[0,1], y\in[1,2]$.
（2）因为$C_2: \begin{cases}x=2\cos\alpha\\y=2\sin\alpha\end{cases}$（$\alpha$为参数，$\frac{\pi}{2}<\alpha<\pi$），整理得$x^2+y^2=4$，表示圆心为$O(0,0)$，半径为2，且位于第二象限的圆弧，如图所示，若直线$y=x+m$过$(1,1)$，则$1=1+m$，解得$m=0$；若直线$y=x+m$，即$x-y+m=0$与$C_2$相切，则$\frac{|m|}{\sqrt{2}}=2$，且$m>0$，解得$m=2\sqrt{2}$，若直线$y=x+m$与$C_1, C_2$均没有公共点，则$m>2\sqrt{2}$或$m<0$，即实数$m$的取值范围$(-\infty,0)\cup(2\sqrt{2},+\infty)$.
\end{solution}
\end{question}

\begin{question}[points=10]
[选修4-5]（10分）

已知$f(x)=|2x|+|x-2|$.

（1）求不等式$f(x)\le6-x$的解集；

（2）在直角坐标系$xOy$中，求不等式组$\begin{cases}f(x)\le y\\x+y-6\le0\end{cases}$所确定的平面区域的面积．
\begin{solution}

\textbf{分析：}（1）分段去绝对值符号求解不等式作答.
（2）作出不等式组表示的平面区域，再求出面积作答.

\textbf{详解：}（1）依题意，$f(x)=\begin{cases}3x-2, & x>2\\x+2, & 0\le x\le2\\-3x+2, & x<0\end{cases}$，不等式$f(x)\le6-x$化为：$\begin{cases}x>2\\3x-2\le6-x\end{cases}$或$\begin{cases}0\le x\le2\\x+2\le6-x\end{cases}$或$\begin{cases}x<0\\-3x+2\le6-x\end{cases}$，解$\begin{cases}x>2\\3x-2\le6-x\end{cases}$得无解；解$\begin{cases}0\le x\le2\\x+2\le6-x\end{cases}$得$0\le x\le2$；解$\begin{cases}x<0\\-3x+2\le6-x\end{cases}$得$-2\le x<0$，因此$-2\le x\le2$，所以原不等式的解集为$[-2,2]$.
（2）作出不等式组$\begin{cases}f(x)\le y\\x+y-6\le0\end{cases}$表示的平面区域，如图中阴影$\triangle ABC$，由$\begin{cases}y=-3x+2\\x+y=6\end{cases}$，解得$A(-2,8)$，由$\begin{cases}y=x+2\\x+y=6\end{cases}$，解得$C(2,4)$，又$B(0,2), D(0,6)$，所以$\triangle ABC$的面积$S_{\triangle ABC}=\frac{1}{2}|BD|\times|x_C-x_A|=\frac{1}{2}\times|6-2|\times|2-(-2)|=8$.
\end{solution}
\end{question}

